\documentclass[11pt]{article}
\usepackage{amsmath,amsfonts,amssymb,amsthm}
\usepackage{geometry}
\usepackage{authblk}
\usepackage{hyperref}
\usepackage{graphicx}
\usepackage{enumitem}
\usepackage{physics}
\usepackage{mathrsfs}
\usepackage{bm}

\geometry{margin=1in}
\title{\textbf{Cross-Absolute Interactions and Force Emergence in the Enhanced Mathematical Ontology of Absolute Nothingness}}

\author[1]{Ricardo Jorge Do Vale}
\author[1]{Advanced Theoretical Physics Institute}
\date{March 21, 2025}

\begin{document}

\maketitle

\begin{abstract}
This paper presents a rigorous formulation of cross-absolute interaction theory within the Enhanced Mathematical Ontology of Absolute Nothingness. We define a generalized interaction framework between absolutes, including transfinite interaction chains, mediated dynamics, force emergence, and information transfer. We demonstrate that all four fundamental forces emerge from compositional interaction operators among orthogonal absolutes. We also define a unified selection principle that governs which interaction configurations manifest in observable physical reality, balancing least action with maximum integrated information. This model extends beyond conventional quantum field theory, embedding all physical interactions in a deeper ontological framework that is both falsifiable and computationally executable.
\end{abstract}

\tableofcontents

\newpage

\section{Introduction and Motivation}

Contemporary physics lacks a consistent model of pre-physical ontology. Space, time, and matter are treated as fundamental rather than emergent. The Enhanced Mathematical Ontology of Absolute Nothingness addresses this foundational problem by introducing \textit{absolutes}: orthogonal, non-physical entities whose interactions generate physical reality.

This paper develops the formal framework for cross-absolute interactions, including the emergence of spacetime curvature, gauge fields, entanglement, and probabilistic phenomena. Our goal is to reconstruct physics from first principles that predate geometry, causality, and measurement.

\section{Absolutes: Ontological Foundations}

\subsection{Definition of Absolutes}

\textbf{Definition 1.} An absolute $\mathcal{A}_i$ is defined by a 4-tuple:
\[
\mathcal{A}_i = \left( \mathcal{S}_i, \mathcal{P}_i, D_i, O \right)
\]
where:
\begin{itemize}[noitemsep]
    \item $\mathcal{S}_i$: structure space (configuration manifold)
    \item $\mathcal{P}_i$: property space (attributes)
    \item $D_i$: internal dynamic generator, $D_i : \mathcal{S}_i \rightarrow \mathcal{S}_i$
    \item $O(\mathcal{A}_i, \mathcal{A}_j) \in [0,1]$: orthogonality metric
\end{itemize}

Absolutes are not spatial, temporal, or material — they exist outside classical constructs. Their interactions create physical constructs such as mass, spin, charge, and curvature.

\subsection{Notation and Conventions}

\begin{itemize}[noitemsep]
    \item $I_{ij}$: interaction operator between $\mathcal{A}_i$ and $\mathcal{A}_j$
    \item $M_{ij}$: mediator space
    \item $E_F$: force extraction operator
    \item $O(\mathcal{A}_i, \mathcal{A}_j)$: orthogonality
    \item $H$: entropy measure
\end{itemize}

\section{Cross-Absolute Interaction Framework}

\subsection{Interaction Operators}

\textbf{Definition 2.} For absolutes $\mathcal{A}_i$, $\mathcal{A}_j$, define:
\[
I_{ij} : \mathcal{S}_i \times \mathcal{S}_j \rightarrow \mathcal{S}_{\text{interaction}}
\]

\textbf{Properties}:
\begin{align*}
I_{ij}(s_i, s_j) &= I_{ji}(s_j, s_i) && \text{(Symmetry)} \\
I_{ii}(s_i, s_i') &= D_i(s_i, s_i') && \text{(Self-interaction)} \\
I_{ij} \circ (I_{jk} \times \text{id}) &= I_{ik} \circ (\text{id} \times I_{ij}) \circ \alpha_{i,j,k} && \text{(Associativity)}
\end{align*}

\textbf{Axiom 1 (Information Conservation)}:
\[
H\left(\bigcup_i \mathcal{S}_i\right) = \sum_i H(\mathcal{S}_i) - \sum_{i<j} H(I_{ij}(\mathcal{S}_i, \mathcal{S}_j))
\]

\subsection{Mediator Spaces and Orthogonality}

\textbf{Definition 3.} Universal structure space:
\[
\mathcal{U} = \bigcup_i \mathcal{S}_i
\]

\textbf{Definition 4.} Mediator space:
\[
M_{ij} = \left\{ m \in \mathcal{U} \; \middle| \; \exists f_i, f_j : \mathcal{S}_i, \mathcal{S}_j \rightarrow m \; \text{structure-preserving} \right\}
\]

\textbf{Theorem (Mediated Interaction)}: Even if $O(\mathcal{A}_i, \mathcal{A}_j) = 1$,
\[
I^{\text{eff}}_{ij}(s_i, s_j) = \min_{m \in M_{ij}} D(f_i(s_i), f_j(s_j))
\]

\subsection{Transfinite Chains}

\textbf{Definition 5.} Transfinite chain of order $\alpha$:
\[
C^\alpha_{ij} = \{ (k_1, ..., k_\beta) \; | \; \beta < \alpha, \; k_1=i, k_\beta=j, \; I_{kl_{l} l_{l+1}} \neq 0 \}
\]

\textbf{Total Interaction Strength}:
\[
S_{\text{total}}(i,j) = \sum_{\alpha < \Omega} \sum_{c \in C^\alpha_{ij}} \left( \prod_{l=1}^{|c|-1} \|I_{c_l c_{l+1}}\| \cdot e^{-\lambda |c|} \right)
\]

\section{Force Emergence Mechanism}

\subsection{Force Extraction}

\textbf{Definition 6.} For each force $F_k$, define:
\[
E_{F_k} \left( \{I_{ij}\} \right) = \sum_{i,j} \alpha^{(k)}_{ij} P_k(I_{ij}) + \sum_{i,j,l} \beta^{(k)}_{ijl} P_k(I_{ij} \circ I_{jl})
\]

\subsection{Force Decomposition}

\[
F_k = \sum_i \gamma^{(k)}_i F_k^{(i)} + \sum_{i<j} \delta^{(k)}_{ij} F_k^{(ij)}
\]

\subsection{Explicit Force Generation}

\begin{align*}
F^{\text{gravity}}_{\mu\nu} &= \sum_{i,j} \alpha^{(g)}_{ij} \frac{\partial^2}{\partial X^\mu \partial X^\nu} \text{Tr}[I_{ij}] \, g^{\mu\nu} \\
F^{\text{EM}}_{\mu\nu} &= \sum_{i,j} \alpha^{(e)}_{ij} \left( \partial_\mu V^{(j)}_\nu - \partial_\nu V^{(i)}_\mu \right)
\end{align*}

\section{Information Transfer Between Absolutes}

\textbf{Definition 7.} Information transfer operator:
\[
T_{ij} : \mathcal{P}_i \rightarrow \mathcal{P}_j
\]

\textbf{Max Transfer Capacity}:
\[
C_{ij} = (1 - O(\mathcal{A}_i, \mathcal{A}_j)) \cdot \min\{H(\mathcal{P}_i), H(\mathcal{P}_j)\}
\]

\textbf{Entanglement Measure}:
\[
\mathcal{E}(\mathcal{A}_i, \mathcal{A}_j) \leq \sqrt{1 - O(\mathcal{A}_i, \mathcal{A}_j)}
\]

\section{Unified Selection Principle}

\textbf{Principle 1 (Least Interaction Action)}:
\[
S_{\text{int}} = \int_\Omega \left( \sum_{i,j} \|I_{ij}\|^2 + \sum_{i,j,k} \lambda_{ijk} \text{Tr}(I_{ij} \circ I_{jk}) \right) d\mu(\Omega)
\]

\textbf{Principle 2 (Maximum Integrated Information)}:
\[
\Phi = H\left( \bigcup_i \mathcal{S}_i \right) - \sum_i H(\mathcal{S}_i)
\]

\textbf{Pareto Selection:}
\[
\text{Realized configuration} = \arg \min S_{\text{int}} \quad \text{and} \quad \arg \max \Phi
\]

\section{Conclusion}

We have established a formal, computable, and ontologically grounded framework in which all physical forces and fields arise from first-order and higher-order interactions between fundamental absolutes. This theory replaces space-time and matter as axioms with emergent consequences of absolute dynamics, making testable predictions and enabling a new class of simulations grounded in transfinite structure spaces and categorical constraints.

\bibliographystyle{plain}
\bibliography{references}

\end{document}
